\documentclass[10pt,a4paper]{article}
\usepackage{amsmath,amssymb,amsthm,graphicx}
\usepackage{hyperref}
\usepackage{latexsym}
\usepackage{amssymb}
\usepackage{amsthm}
\usepackage{ulem}
\usepackage{enumitem}
\setlist{itemsep=0.4pt,topsep=0pt}
\parindent 0pt
\parindent 0pt
\parskip 12pt
\newcommand{\HCF}{\mathrm{hcf}}
\newcommand{\ZZ}{\mathbb{Z}}
\newcommand{\rk}{\mathrm{rk}}
\DeclareMathOperator{\len}{length}
\DeclareMathOperator{\bl}{b}
\DeclareMathOperator{\cl}{cl}
\DeclareMathOperator{\di}{d}
\DeclareMathOperator{\Cl}{Cl}
\DeclareMathOperator{\acl}{acl}
\DeclareMathOperator{\tp}{tp}
\DeclareMathOperator{\Th}{Th}
\newcommand{\Cor}{\mathcal{C}_\mathrm{or}}
\newcommand{\Cun}{\mathcal{C}_\mathrm{un}}
\newcommand{\mons}{\mathbb{M}}
\newcommand{\strong}{\mathrm{strong}}
\newcommand{\axioms}{\Delta_{\mathrm{FSTS}}}
\newtheoremstyle{note}% hnamei
{8pt}% hSpace abovei
{14pt}% hSpace belowi
{\normalfont}% hBody fonti
{}% hIndent amounti1
{\bfseries}% hTheorem head fonti
{ }% hPunctuation after theorem headi
{.5em}% hSpace after theorem headi2
{}% hTheorem head spec (can be left empty, meaning ‘normal’)i
\theoremstyle{note}
\newtheorem{thm}{Theorem}[section]
%\newtheorem{theorem}{Theorem}
\newtheorem{lemma}[thm]{Lemma} 
\newtheorem{df}[thm]{Definition}
\newtheorem{rmk}[thm]{Remark} 
\newtheorem{prop}[thm]{Proposition}
\newtheorem{fact}[thm]{Fact}
\newtheorem{cor}[thm]{Corollary}
\newtheorem{conj}[thm]{Conjecture}
\newtheorem{ex}[thm]{Example}
\newtheorem{qu}[thm]{Question}
\newtheorem*{notation}{Notation} 
\title{A set closed under algebraic closure is an elementary subsystem}
\author{Silvia Barbina, Enrique Casanovas}
\date{arxiv:2411.13723v3; Creative Commons Attribution 4.0 International}
\begin{document}
\maketitle
\noindent\textit{Editorial note.} The counted statement is the author's environment at source0.tex lines 789-790 and its proof environment opens at line 791; the pairing is by adjacency and the retained proof range is the whole environment.\par\medskip
\noindent\textit{Editorial note.} The preamble is the author's own package lines, macro block and theorem-environment declarations (source lines 4-65); the author's optioned amsmath line and the fleqn-dependent layout line are not retained.\par\medskip
\noindent\textit{Editorial note.} The source carries no label-number dump and no counter values are imposed, so the printed numbering is the wrapper's default; identity is the author label, or the pinned source position where the author gave none.\par\medskip
\noindent\textit{Editorial note.} Metadata provenance: the source record of the pinned copy and the licence evidence file are both present; the licence is CC BY 4.0 with status supported, and the card's source\_version matches the value recorded in that evidence file.\par\medskip
\noindent\textit{Editorial note.} r002 note: r001 failed because the card's own materials prose, which the wrapper typesets, carried characters the font does not provide; r002 replaces them by equivalent TeX commands and changes nothing else, and no retained author line is altered.\par\medskip
\noindent\textit{Editorial note.} r003 note: the retained author proof cites the bibliography keys pmneumann and hodges, so r003 bundles the article's own bibliography file; r002 had been built with the wrong flag and changed nothing.\par\medskip
\noindent\textit{Editorial note.} r004 note: the retained proof cites the bibliography keys pmneumann and hodges, which the prepared per-source bibliography of this source does not carry; r004 adds a task-directory bibliography.tex holding those two entries copied character for character from the article's own thebibliography block and points bibliography\_paths at it. No bibliographic datum is invented and no retained author line is altered.\par\medskip
\noindent\textit{Editorial scope.} One original result of the pinned author manuscript, counted once: a set equal to its algebraic closure is an elementary subsystem (author label Feb\_23\_2024\_8).\par\medskip
\noindent\textit{Allowed standard prerequisites (editorial).}\begin{itemize}
\item Algebraic closure. The algebraic closure of A is the set of elements whose type over A is algebraic, and it is itself closed under the closure.. It is the hypothesis A equals its algebraic closure.
\item Unconfined system. A Steiner triple system admitting an HF-ordering without maximum.. It is the ambient hypothesis.
\item Neumann's lemma. Given elements outside the algebraic closure of A, any set in the monster model can be moved so that the chosen elements avoid it, over A.. It supplies the back-and-forth move used in the proof.
\end{itemize}\medskip
\begin{lemma}\label{Feb_23_2024_7} If $<$ is an HF-ordering of $B$  and $A=\cl_<(A)\subseteq B$, then  $<$ is an HF-ordering of $B$ over $A$. 
\end{lemma}

	\begin{prop} \label{Feb_23_2024_8} If  $M$ is an unconfined STS  and $A=\acl(A)\subseteq M$, then  $A\leq M$.
	\end{prop}

	\begin{proof}   We can expand the monster model $\mons$ by adding an HF-ordering $<$ of $\mons$, but we work in $\mons$.   
		By P.M. Neumann's Lemma (2.3 in \cite{pmneumann} --- also \cite[Corollary~4.2.2]{hodges}),  given  $a_1,\ldots,a_n\not\in \acl(A)$, for any set $B$ in $\mons$  there are $b_1\ldots b_n$ with $b_1,\ldots,b_n\equiv_A a_1\ldots a_n$ such that  $b_1,\ldots,b_n\not\in B$. Hence if $M\smallsetminus A=\{a_i\mid i<\kappa\}$  and  $p(x_i\mid i<\kappa)=\tp((a_i\mid i<\kappa)/A)$, then,  for any set $B$,  the type~$p$ has a realization  $(b_i\mid i<\kappa) $ such that   $\{b_i\mid i<\kappa\}\cap B=\emptyset$. Using this, for any $\lambda$  we can construct a sequence of models $(M_i\mid i<\lambda)$  containing $A$ such that  $M=M_0$, $M_i\cong_A M_j$  and $M_i\cap M_j=A$  for every $i\neq j$. Let $B=\cl_<(A)\subseteq \mons$ and let  $\mu= |B|= |A|+\omega$.  We may assume $\mu <\lambda$, which implies  $M_i\cap B = A$  for some $i$ (if $A\subsetneq M_i\cap B$ for every $i<\lambda$, we can choose $a_i\in (M_i\cap B)\smallsetminus A$ and the $a_i$ are all distinct, which contradicts $|B|<\lambda$).  By Lemma~\ref{Feb_23_2024_7}, $<$ is an HF-ordering over $B$. Therefore $<_i=<\restriction M_i$ is an HF-ordering of $M_i$ over  $A= B\cap M_i$. Since $M_i\cong_A M$,  there is an HF-ordering of $M$ over $A$.
	\end{proof}

% Verbatim extract of the two bibliography entries cited by the retained author proof of this card,
% copied character for character from the pinned source (its own thebibliography block).
\begin{thebibliography}{10}
\bibitem{pmneumann}
P.M. Neumann.
\newblock The structure of finitary permutation groups.
\newblock {\em Archiv der Mathematik}, 27:3--17, 1976.
\bibitem{hodges}
W.~Hodges.
\newblock {\em Model theory}.
\newblock Cambridge University Press, 1993.
\end{thebibliography}
\end{document}
